\chapter{Asbestosis}

\section*{Definition and Overview}
Asbestosis is a form of interstitial pulmonary fibrosis caused by the inhalation of asbestos fibres over a prolonged period. The inhaled fibres lodge deep within the lungs and provoke a chronic inflammatory response that progressively replaces normal air-exchange tissue with scar tissue, making the lungs stiff and difficult to expand. It is an occupational lung disease that typically becomes symptomatic many years after exposure. Asbestosis is distinct from other asbestos-related conditions, including benign pleural plaques, diffuse pleural thickening, lung cancer, and mesothelioma, although several of these may coexist in the same person.

\section*{Epidemiology}
Asbestosis occurs almost exclusively in people who have worked with asbestos, most commonly in industries such as mining, milling, shipbuilding, construction, insulation, and brake and boiler maintenance. Because symptoms appear only after a long latency, usually 15 to 30 years or more, cases continue to be diagnosed even though most countries now ban or strictly regulate the use of asbestos. Men are affected far more often than women, reflecting occupational exposure patterns, and the number of new cases has fallen as workplace protections have improved and asbestos use has been phased out. Secondary exposure in household members of asbestos workers has also been reported.

\section*{Aetiology and Pathophysiology}
\begin{itemize}
    \item \textbf{Cause:} prolonged inhalation of asbestos fibres small enough to reach the alveoli, the deepest part of the lungs
    \item \textbf{Mechanism:} the fibres cannot be cleared effectively and irritate the lung tissue, attracting immune cells that release inflammatory and fibrogenic chemicals
    \item \textbf{Fibrosis:} repeated irritation leads to progressive scarring of the delicate gas-exchange tissue, reducing the transfer of oxygen into the blood
    \item \textbf{Risk factors:} higher fibre concentrations and longer duration of exposure increase the risk; cigarette smoking greatly increases the risk of asbestos-related lung cancer but is less central to asbestosis itself
    \item \textbf{Related effects:} pleural plaques, diffuse pleural thickening, and pleural effusions may develop separately from the pulmonary fibrosis
\end{itemize}

\section*{Clinical Features}
\begin{itemize}
    \item Gradually increasing breathlessness, first noticed on exertion and progressing over years
    \item A persistent, usually dry cough
    \item Chest tightness or discomfort
    \item Fine crackling sounds (crepitations) at the lung bases heard with a stethoscope
    \item Finger clubbing in some cases
    \item Reduced exercise tolerance and fatigue
    \item Symptoms typically appear many years after the original exposure
\end{itemize}

\section*{Differential Diagnosis}
\begin{itemize}
    \item Chronic obstructive pulmonary disease (COPD) and emphysema, especially in smokers
    \item Idiopathic pulmonary fibrosis and other interstitial lung diseases
    \item Silicosis and other pneumoconioses caused by inhaled dust
    \item Heart failure causing pulmonary oedema
    \item Tuberculosis or other chronic infections of the lung
    \item Lung cancer or mesothelioma, which may coexist with asbestosis
\end{itemize}

\section*{When to Seek Medical Care}
See a doctor if you have ever worked with asbestos and develop breathlessness, a persistent cough, or chest symptoms, and tell the doctor about your work history. Anyone with significant past exposure should also discuss regular monitoring with their health-care team.

\textbf{Call 000 (or attend an emergency department) for:}
\begin{itemize}
    \item Sudden severe breathlessness or difficulty breathing
    \item Coughing up blood
    \item Sudden or crushing chest pain
    \item Collapse, fainting, or blue lips or face
\end{itemize}

\section*{Diagnosis and Investigations}
\begin{itemize}
    \item \textbf{History:} a detailed occupational and environmental exposure history is the single most important diagnostic clue
    \item \textbf{Examination:} auscultation of the lungs and inspection of the fingers for clubbing
    \item \textbf{Chest X-ray:} may show small irregular shadows, particularly in the lower zones, together with pleural changes
    \item \textbf{High-resolution CT scan:} the most sensitive imaging for detecting fibrosis, honeycombing, and pleural plaques
    \item \textbf{Lung function tests (spirometry):} characteristically show a restrictive pattern with reduced lung volumes and a reduced gas transfer factor
    \item \textbf{Blood tests:} may help to exclude autoimmune and other causes of pulmonary fibrosis
    \item \textbf{Bronchoscopy or lung biopsy:} occasionally required when the diagnosis remains uncertain
\end{itemize}

\section*{Management}
\begin{itemize}
    \item \textbf{No cure:} treatment focuses on slowing progression and controlling symptoms
    \item \textbf{Stopping further exposure:} avoid any continuing contact with asbestos
    \item \textbf{Smoking cessation:} essential, because smoking markedly increases the risk of lung cancer in asbestos-exposed people
    \item \textbf{Pulmonary rehabilitation:} supervised exercise and education to improve fitness, breathlessness, and daily function
    \item \textbf{Oxygen therapy:} provided for people with low blood oxygen levels, particularly during activity or at night
    \item \textbf{Medications:} treatment of cough, respiratory infection, and cardiac complications as they arise, and antifibrotic therapy in selected cases
    \item \textbf{Vaccinations:} annual influenza and pneumococcal vaccines to reduce the risk of serious respiratory infection
    \item \textbf{Lung transplantation:} considered for eligible patients with advanced disease
\end{itemize}

\section*{Complications}
\begin{itemize}
    \item Progressive breathlessness, disability, and respiratory failure
    \item Low oxygen levels in the blood (hypoxaemia)
    \item Pulmonary hypertension and strain on the right side of the heart (cor pulmonale)
    \item Recurrent respiratory infections
    \item Lung cancer, which is far more likely in exposed people who smoke
    \item Mesothelioma, a cancer of the pleural lining, which is a separate risk of asbestos exposure
\end{itemize}

\section*{Prognosis and Outcomes}
Asbestosis is a slowly progressive disease, but the rate of deterioration varies greatly between individuals. It may remain stable for long periods in some people, while others gradually develop increasing breathlessness and disability over many years. Progression often slows once exposure has stopped. The prognosis is worse in people who continue to smoke or who develop lung cancer, mesothelioma, or cardiac complications. Many people continue to live active, productive lives for a long time after diagnosis, particularly with good supportive care and rehabilitation.

\section*{Prevention and Self-Care}
\begin{itemize}
    \item The most effective prevention is to avoid inhaling asbestos fibres
    \item Follow workplace safety rules, and use approved respirators and protective clothing wherever asbestos is handled
    \item Never disturb, cut, sand, or drill asbestos materials in the home; use licensed asbestos removal services
    \item Do not smoke, and seek help to quit if needed
    \item Keep up to date with influenza and pneumococcal vaccinations
    \item If previously exposed, attend regular check-ups and lung function monitoring, and report new symptoms early
\end{itemize}

\section*{Sources and Further Reading}
The following organisations provide reliable, evidence-based information on asbestosis for patients and healthcare students.

\begin{itemize}
    \item NHS. \textit{Asbestosis: symptoms and treatment.} https://www.nhs.uk/conditions/asbestosis/
    \item National Heart, Lung, and Blood Institute (NHLBI). \textit{Asbestosis information.} https://www.nhlbi.nih.gov/health/asbestosis
    \item Centers for Disease Control and Prevention (CDC). \textit{Asbestos and health.} https://www.cdc.gov/asbestos/
    \item Occupational Safety and Health Administration (OSHA). \textit{Asbestos and worker safety.} https://www.osha.gov/asbestos
    \item Mayo Clinic. \textit{Asbestosis overview.} https://www.mayoclinic.org/diseases-conditions/asbestosis/
    \item MedlinePlus. \textit{Asbestosis.} https://medlineplus.gov/asbestosis.html
\end{itemize}
